\section{Belastung der einzelnen Stäbe}
\subsection{Stab 5}
Da Stab 5 belastungen nur an seiner Endpunkte hat, sind die 2 Kräfte in diesen Punkte im Stabrichtung gerichet, und sind entgegengesetzt gleich groß.
\begin{figure}[H]
  \centering
  \begin{tikzpicture}[thick]
    \draw (C) -- (A);
    \draw[force] (C) -- ($ (C)!-1cm!(A) $) node[above] {$S_5$};
    \draw[force] (A) -- ($ (A)!-1cm!(C) $) node[above] {$S_5$};
    \draw[force, color=gray!50] (A) -- ++(1, 0) node[right] {$A_X$};
    \draw[force, color=gray!50] (A) -- ++(0, 1) node[right] {$A_Y$};
  \end{tikzpicture}
  \caption{Belastung von Stab 5}
\end{figure}
\begin{align*}
  |S_5| = |\mathbf{A}| = \sqrt{A_X^2 + A_Y^2} = 16,802 \, \text{kN}
\end{align*}
Da $A_X$ und $A_Y$ auch negativ sind:
\begin{align*}
  \boxed{
    S_5 = 16,802 \, \text{kN}
  }
\end{align*}

\subsection{Stab 4}
Nach dem selben Logik als Stab 5:
\begin{align*}
  &S_4 = G \\
  &\boxed{
    S_4 = 3 \, \text{kN}
  }
\end{align*}

\subsection{Stab 2 und 3}
Im Gelenk B sind 3 Kräfte zu finden. Die Summe der 3 Kräften muss 0 sein, für die Konstruktion im Gleichgewicht zu sein.
\begin{figure}[H]
  \centering
  \begin{tikzpicture}[thick]
    \draw[force, color=gray!50] (B) -- ++(1, 0) node[right] {$B_X$};
    \draw[force, color=gray!50] (B) -- ++(0, 1) node[above] {$B_Y$};
    \draw[dashed, color=gray!50] ($(B) + (1, 0)$) -- ++(0, 1);
    \draw[dashed, color=gray!50] ($(B) + (0, 1)$) -- ++(1, 0);
    \draw[force] (B) -- ++(1, 1) node[above right] {$\mathbf{B}$};

    \draw[force] (B) -- ++(-60:1.4) node[right] {$\mathbf{S_3}$};
    \draw[force] (B) -- ++(-120:1.4) node[left] {$\mathbf{S_2}$};
    \hinge {B};
  \end{tikzpicture}
  \caption{Belastung von Gelenk B}
\end{figure}

\setcounter{equation}{0}
\begin{align}
  \sum F_X = B_X + S_3 \cdot cos(-60^\circ) + S_2 \cdot cos(-120^\circ) = 0\\
  \sum F_Y = B_Y + S_3 \cdot sin(-60^\circ) + S_2 \cdot sin(-120^\circ) = 0
\end{align}
\begin{align*}
  \boxed{
    \begin{aligned}
      S_2 &= -20{,}928 \, \text{kN}
  \end{aligned}}
  &\qquad
  \boxed{
    \begin{aligned}
      S_3 &= 11{,}916 \, \text{kN}
  \end{aligned}}
\end{align*}
